\documentclass{article}
\usepackage{tikz}
\usetikzlibrary{arrows.meta, positioning, shapes.geometric, calc}
\usepackage{graphicx}
\usepackage{url}
\usepackage{booktabs}
\usepackage{tabularx}
\usepackage{multirow}
\usepackage{array}
\usepackage{amsmath}
\usepackage{subcaption}
\usepackage{float}
\usepackage{textcomp}
\usepackage[version=4]{mhchem}
\usepackage[a4paper,margin=2.5cm]{geometry}
\usepackage{eurosym}
\usepackage{threeparttable}
\usepackage{xcolor}  % for \color in title rule

\title{\textbf{Andorra Digital Twin: A Data-Grounded Scenario Engine for National Development Planning}\\[0.1em]
{\normalfont\rule{0.6\textwidth}{0.4pt}}}

\date{15th December 2025}
\author{}

\begin{document}

\maketitle

\begin{center}
\small
\begin{tabular}{@{}>{\centering\arraybackslash}p{0.44\textwidth}@{\hspace{0.08\textwidth}}>{\centering\arraybackslash}p{0.44\textwidth}@{}}
  \textbf{Marcel Bartumeu Ramentol}\newline
  MIT Media Lab\newline
  Massachusetts Institute of Technology\newline
  Cambridge, MA, USA
  &
  \textbf{Parfait Atchade-Adelomou}\newline
  MIT Media Lab\newline
  Massachusetts Institute of Technology\newline
  Cambridge, MA, USA
  \\[1.2em]
  \textbf{Luis Alonso-Pastor}\newline
  MIT Media Lab\newline
  Massachusetts Institute of Technology\newline
  Cambridge, MA, USA
  &
  \textbf{Marta Domenech}\newline
  Andorra Research + Innovation (AR+I)\newline
  Andorra la Vella, Andorra
  \\[1.2em]
  \textbf{Adrian Mora-Carrero}\newline
  Escuela T\'ecnica Superior de Ingenieros Inform\'aticos\newline
  Universidad Polit\'ecnica de Madrid\newline
  Madrid, Spain
  &
  \textbf{Kent Larson}\newline
  MIT Media Lab\newline
  Massachusetts Institute of Technology\newline
  Cambridge, MA, USA
\end{tabular}
\end{center}

\vspace{1.5em}

% ============================================================
\section{Abstract}
% ============================================================

Planning future national development under uncertainty is difficult, especially
for small countries where data are fragmented and frequently updated. While
digital twins and scenario planning can support decision-making without relying
on a single forecast, they often lack historical grounding, cross-domain
consistency, and reproducible update pipelines. We present the Andorra Clone
Scenario Engine, a data-grounded framework that harmonizes multi-source national
indicators and produces structured future scenarios for Andorra under four
contrasting pathways: \textit{Continuity}, \textit{Overgrowth},
\textit{Degrowth}, and \textit{Density}, with a planning horizon of 2049.
A central methodological contribution is the inversion of the population--economy
causal chain: rather than treating population as an exogenous driver, the model
derives population endogenously from GDP growth via a labour-immigration
elasticity (0.40 pp population growth per 1 pp GDP growth at a one-year lag,
0.10 pp at two years), grounded in IMF\,(2025) and World Bank data for Andorra.
Three core parameters---total factor productivity growth, the tourism--GDP
elasticity, and the household-income coefficient---are calibrated against
official national statistics for 2010--2024 via Nelder--Mead optimisation
(RMSE\,=\,1.77\% on real GDP per capita), anchoring projections in observed
dynamics rather than prior assumptions alone.
Two structural nonlinearities are retained: a housing-affordability feedback
that suppresses net in-migration when the rent-to-income ratio exceeds 30\%, and
a logistic tourism carrying-capacity ceiling grounded in the Butler\,(1980) TALC
framework. Parametric uncertainty
is quantified through 500 Monte Carlo draws, yielding P10--P90 bands across all
output variables. By 2049, projected populations range from 73{,}172 (Degrowth)
to 179{,}021 (Overgrowth), with the Continuity baseline at 117{,}666 and the
Density compact-growth pathway at 104{,}914. The resulting scenarios support
long-term infrastructure planning, sustainability assessment, and policy
comparison, and provide a baseline for future human-in-the-loop and agent-based
simulations. In line with the MIT Media Lab City Science vision, the framework
is designed as a transparent, human-in-the-loop decision-support baseline for
exploring trade-offs across plausible futures. Limitations include residual data
gaps, uncertainty in exogenous drivers, simplified causal assumptions, and the
absence of endogenous climate forcing, but the approach offers a reproducible,
updateable, and historically grounded foundation for scenario baselines in
data-scarce contexts.

% ============================================================
\section{Introduction}
% ============================================================

Small countries increasingly need evidence-based ways to explore long-term
planning choices under climate pressure, resource constraints, and tightly
coupled social--economic--environmental systems~\cite{riahi2017ssps}. Yet many
planning workflows still rely on single ``most-likely'' projections, which can
miss uncertainty, nonlinear effects, and cross-sector
dependencies~\cite{amer2013scenario}. Although machine learning approaches are
widely used for forecasting in urban and territorial
systems~\cite{batty2018digitaltwins}, they often require large, stable datasets
and can be hard to interpret---constraints that are especially acute in
small-country contexts where indicators are fragmented across institutions and
across time.

Digital twins and scenario-based decision support offer a promising alternative
by organizing heterogeneous evidence into interactive models that compare
multiple plausible futures rather than one deterministic
outcome~\cite{batty2018digitaltwins}. However, existing frameworks often
struggle with transparent assumptions, consistent indicator definitions, and
reproducible pipelines that remain updateable as new official statistics are
released~\cite{udtreview2024}. Without a standardized scenario baseline,
scenario narratives can become incomparable, and trajectories across domains
(e.g., population, housing, energy, land) can drift into internal
inconsistency. This challenge is especially relevant to the MIT Media Lab City
Science vision, which emphasizes transparent, human-centered, and
decision-oriented tools that help communities explore trade-offs across
alternative futures rather than optimize a single
forecast~\cite{alonso2018cityscope,noyman2022cityscope}.

This paper introduces the Andorra Clone Scenario Engine, a data-grounded
framework that harmonizes multi-source national indicators into a reproducible
KPI dataset and generates structured scenario trajectories for Andorra to 2049.
A central innovation relative to earlier versions of the framework is a
\textbf{methodological inversion of the population--economy causal chain}.
Whereas a population-as-driver approach treats future headcount as an
exogenous input and then derives economic and infrastructure consequences,
Andorra's demographic structure---where roughly 55\% of residents are
foreign-born labour migrants whose presence directly follows economic
activity---motivates the reverse direction: GDP growth is the primary
exogenous lever, and population evolves endogenously through a lagged
labour-immigration elasticity grounded in IMF\,(2025) Andorra Selected Issues
and World Bank panel data. A KPI dictionary and data pipeline standardize units,
frequency, and provenance, while the scenario engine encodes explicit drivers
and propagates them through bounded, coupled update rules to maintain
cross-domain consistency. Key state variables such as building stock evolve
endogenously via a permits--demolition stock--flow mechanism, and downstream
infrastructure and environmental indicators are computed deterministically from
the shared state. The result
is a transparent, updateable, and auditable scenario baseline suitable for
infrastructure stress testing, sustainability assessment, and policy comparison
in data-scarce environments, and it provides a foundation for future extensions
that incorporate humanized AI agents to study behavioral responses under
alternative futures.

% ============================================================
\section{Related Work}
% ============================================================

Scenario engines and digital twins for planning sit at the intersection of
structured futures, territorial-scale decision support, and reproducible
indicator pipelines. This section summarizes key strands and positions the gap
addressed by the Andorra Clone Scenario Engine.

\subsection{Scenario Planning and Structured Futures}
Scenario planning explores plausible futures under deep uncertainty by
constructing internally coherent alternative trajectories for stress testing
rather than point prediction~\cite{amer2013scenario}. In climate research,
scenario families such as the Shared Socioeconomic Pathways (SSPs) provide
narratives and quantitative pathways to support consistent cross-domain
analysis~\cite{riahi2017ssps,oneill2014ssp}. However, SSP-style approaches are
not tailored to small-country workflows where indicators are fragmented, time
series are short, and update rules must remain interpretable as new official
statistics arrive. This motivates lighter scenario engines that preserve
explicit assumptions, comparability, and auditability while adapting to
data-scarce contexts.

\subsection{Digital Twins for Territorial Planning}
Digital twins connect data, models, and decision workflows through maintainable
digital representations of real systems, where value comes from purposeful
abstraction and interpretability rather than perfect
mirroring~\cite{batty2018digitaltwins}. Reviews emphasize recurring challenges
in data integration, interoperability, governance, and auditable
assumptions~\cite{udtreview2024}. National initiatives likewise stress
trustworthy data and clear evolution
mechanisms~\cite{cdbb2018gemini,ndtp2024principles}. Yet many programs focus on
platform architecture and leave open the practical problem of building an
\emph{updateable, reproducible scenario baseline} that stays cross-domain
consistent (population--housing--economy--environment). The Andorra Clone
Scenario Engine is positioned as an upstream layer that prioritizes harmonized
baselines and consistent scenario propagation before higher-fidelity simulation
or interfaces. This orientation is also consistent with the MIT Media Lab City
Science approach, in which digital platforms such as CityScope are used not as
exhaustive replicas of reality, but as transparent, interactive
decision-support systems for comparing alternative urban futures and supporting
community engagement~\cite{grignard2018cityscope,berke2022using}.

\subsection{Labour-Migration--Economy Linkages}
In small economies with high foreign-born labour shares, the direction of
causality between economic growth and population growth runs primarily from
economy to people rather than vice versa. A growing empirical literature
documents GDP-to-migration elasticities in the range of 0.3--0.6 for small
European destinations~\cite{imf2025andorra,worldbank2024}. Housing costs
introduce a self-limiting feedback: when rent burdens exceed roughly 30\% of
household income, net in-migration is empirically shown to decelerate even
under continued economic expansion~\cite{Sa2015,Diamond2016}. Incorporating
these two mechanisms---a lagged GDP elasticity and an affordability
brake---makes the population projection endogenous and removes the need to
specify future headcounts as arbitrary scenario assumptions.

\subsection{Indicator Frameworks, Harmonization, and Reproducibility}
Frameworks such as OECD Regional Well-Being and the UN SDG indicators provide
guidance on what to measure and how to maintain comparability, including
standardized metadata and evolving
methodologies~\cite{oecdregionalwellbeing,unstatsSDGmetadata}. In small-country
settings, operational use typically requires a pragmatic KPI dictionary that
harmonizes heterogeneous sources into consistent units, frequency, and
definitions. A key challenge is \emph{computational linkability}: ensuring
indicators evolve together under scenario drivers without internal
contradictions. Reproducibility principles (e.g., FAIR) further imply that
provenance, transformations, and definitions should be versioned and re-runnable
as data updates~\cite{wilkinson2016fair}. In practice, ad hoc preprocessing and
non-versioned assumptions remain common, limiting auditability, comparability
across releases, and reuse.

\subsection{Built Environment Dynamics and Environmental Consistency}
Housing affordability is often tracked through cost-burden metrics, while built
environment evolution is frequently represented with stock--flow logic
(construction/permits versus depreciation), responding to demand and
constraints~\cite{eurostatHousingCostOverburden}. In land-constrained settings,
coupling building expansion to land and natural coverage is particularly
important. Environmental trajectories (CO$_2$, air quality, renewables, resource
demand) are commonly linked to socio-economic drivers in integrated scenario
frameworks, but maintaining consistency at national scale is difficult when
indicators are treated as independent trends~\cite{riahi2017ssps}. This
motivates lightweight accounting-style couplings that preserve interpretability
for stress testing while avoiding the false precision of over-specified models in
data-limited settings.

\subsection{Gap Analysis and Contribution of This Work}
Across the literature, three gaps recur when moving from general
scenario/digital twin concepts to deployable small-country workflows: (i)
fragmented baselines that are difficult to harmonize and update, (ii)
cross-domain inconsistency without explicit couplings, and (iii) limited
auditability due to non-versioned transformations and assumptions. The Andorra
Clone Scenario Engine addresses these gaps with a minimal, transparent,
updateable approach: a KPI dictionary and harmonization pipeline that yields a
consistent national state vector; an inverted causal chain that derives
population endogenously from GDP growth rather than imposing it exogenously;
and a scenario engine that produces comparable trajectories via
explicit drivers and bounded update rules. The output is an auditable scenario
baseline that can serve as an upstream layer for broader digital twin interfaces
and later human-in-the-loop or agent-based extensions.

% ============================================================
\section{Methodology}
\label{sec:methodology}
% ============================================================

This section describes the Andorra Clone Scenario Engine as a reproducible
modeling pipeline: a harmonized KPI state representation, a versioned,
updateable data layer, a modular scenario layer (drivers and constraints), a
deterministic yearly state-update engine, a spatial growth allocation layer,
and a derived-infrastructure calculator. The goal is to produce interpretable,
comparable KPI trajectories under alternative assumptions while keeping
provenance and update rules explicit.

\subsection{System Overview and Architecture}
\label{subsec:system_overview}

As shown in Figure~\ref{fig:system_architecture}, the system is structured as a
layered pipeline: (i) a \textit{data layer} that ingests, cleans, and
harmonizes multi-source datasets into a historical baseline state; (ii) an
\textit{engine + calculator layer} that defines scenario and driver
configurations (horizon, exogenous trajectories, policy regimes/constraints)
and advances the yearly state vector through the core scenario engine; (iii) a
\textit{metric layer} that exports scenario outputs (time series, endpoints,
rollups) and computes derived metrics and QA signals (deltas vs.\ baseline,
constraint/feasibility flags, and consistency checks); (iv) a \textit{visualisation layer} that renders
outputs through a lightweight web dashboard and a 3D Unity-based map interface;
and (v) a \textit{tangible interface layer} that enables human-in-the-loop
exploration via an Arduino-controlled physical table for interactive scenario
comparison. This layered architecture is also aligned with the MIT Media Lab
City Science approach, in which harmonized data, simulation, visualization, and
tangible interfaces are integrated to support transparent comparison of
alternative futures rather than a single optimized output~\cite{noyman2019reversed}.

\begin{figure}[H]
    \centering
    \includegraphics[width=0.5\linewidth]{System Architecture.png}
    \caption{End-to-end methodology. The engine converts heterogeneous datasets
    into a harmonized KPI baseline, applies scenario drivers and constraints,
    advances the national state vector yearly via deterministic update rules,
    and renders the results through a web dashboard and a tangible physical interface.}
    \label{fig:system_architecture}
\end{figure}

\subsection{KPI State Representation}
\label{subsec:kpi_framework}

The national system is represented as a yearly KPI state vector
\begin{equation}
S_t = \{k^{(1)}_t, k^{(2)}_t, \dots, k^{(n)}_t\},
\end{equation}
where KPIs are grouped into (i) Social, (ii) Economic, (iii) Built Environment,
(iv) Infrastructure, and (v) Environmental layers. Each KPI is either
\textbf{observed} (initialized from data at baseline year) or \textbf{derived}
(computed from other KPIs through deterministic transformations).

The baseline year state $S_0$ (calibrated to 2024) initializes all runs.
Scenarios do not change the KPI definitions; they change only the
\textit{drivers}, \textit{constraints}, or selected \textit{parameters} used
by the update rules, ensuring that outputs remain comparable across runs and
across future updates of the baseline.

\subsection{Data Pipeline and Harmonization}
\label{subsec:data_pipeline}

The pipeline converts heterogeneous indicator sources into a single harmonized
baseline state and a consistent historical time series used for calibration and
validation. Sources may differ in units, reporting frequency, and missingness;
therefore, the pipeline enforces a consistent yearly timestep and a standard
KPI dictionary.

\paragraph{Transformations and alignment.}
\begin{itemize}
    \item \textbf{Unit normalization:} monetary values are expressed in
          consistent currency and time basis; physical quantities are
          standardized (e.g., kWh, m$^3$).
    \item \textbf{Temporal alignment:} all KPIs are mapped onto a yearly index;
          sub-annual series are aggregated (e.g., sum or mean) using
          KPI-appropriate rules.
    \item \textbf{Missing values:} missing values inside the historical range
          are filled by conservative interpolation; missing values at boundaries
          use last-observation carried forward/backward to preserve continuity
          and should be interpreted as pragmatic imputations rather than
          observations.
    \item \textbf{Bounds and sanity checks:} physically or logically implausible
          values are prevented using conservative clamping ranges to keep
          trajectories interpretable and stable.
\end{itemize}

These steps prioritize consistency and reproducibility, although they do not
remove the need for later source-specific validation.

\subsection{Historical Calibration}
\label{sec:calibration}

Three parameters govern the medium-run level and growth trajectory of GDP per
capita and household income: the total factor productivity growth rate
$g_\text{TFP}$, the tourism--GDP elasticity $\gamma_2$, and the
household-income coefficient $a_1$. Default prior values derived from the
literature ($g_\text{TFP}^0=1.6\%$, $\gamma_2^0=0.05$, $a_1^0=1.296$) are
updated by minimising the root-mean-square error between simulated and observed
real GDP per capita over the period 2010--2024.

Formally, let $\hat{Y}_t(\boldsymbol{\theta})$ denote the simulated GDP per
capita in year $t$ given parameter vector
$\boldsymbol{\theta}=(g_\text{TFP},\gamma_2,a_1)$, and let $Y_t^*$ denote the
observed value. The calibration objective is
\begin{equation}
  \hat{\boldsymbol{\theta}} = \arg\min_{\boldsymbol{\theta}}
    \sqrt{\frac{1}{T}\sum_{t=2010}^{2024}
    \left(\hat{Y}_t(\boldsymbol{\theta}) - Y_t^*\right)^2}.
  \label{eq:calib_obj}
\end{equation}
Equation~\eqref{eq:calib_obj} is solved with the derivative-free Nelder--Mead
simplex algorithm~\cite{Nelder1965} as implemented in
\texttt{scipy.optimize}~\cite{Virtanen2020}.

The calibrated estimates are
\begin{equation}
  \hat{g}_\text{TFP} = 0.865\%~\text{yr}^{-1}, \quad
  \hat{\gamma}_2 = 0.0651, \quad
  \hat{a}_1 = 1.386,
  \label{eq:calib_result}
\end{equation}
yielding an in-sample RMSE of 1.77\%. The lower calibrated $g_\text{TFP}$
relative to the prior (1.60\%) reflects Andorra's post-2008 productivity
slowdown documented by OECD\,(2022); the upward revision of $a_1$ captures the
sustained compression of the wage-to-income gap driven by tourism-sector
earnings. All scenario projections use these calibrated values as their base.

\paragraph{Out-of-target consistency checks.}
Calibration targets GDP per capita exclusively. Two qualitative checks provide
directional reassurance. First, the observed Andorran population was nearly flat
from 2010 to 2013 ($-532$ persons over three years), a period of suppressed net
in-migration consistent with the housing-affordability brake described in
Section~\ref{sec:housing_feedback}; from 2015 onward the population accelerated
(+15{,}365 persons to 2024) as GDP growth eased affordability pressures. The
calibrated model reproduces this qualitative two-regime pattern without
additional tuning. Second, the observed average monthly salary grew from
\EUR\,1{,}995 (2010) to \EUR\,2{,}571 (2024), a compounded annual rate of
approximately 1.8\%; the calibrated $\hat{a}_1=1.386$ coefficient implies a
wage-to-GDPpc ratio for 2010--2024 that tracks this growth profile to within
$\sim$4\% mean absolute deviation.

\subsection{GDP-Driven Endogenous Population Model}
\label{sec:gdp_pop_model}

A structurally important departure from earlier versions of the framework is the
treatment of population as an \textbf{endogenous} variable driven by GDP growth
rather than as an exogenous scenario input. This inversion is motivated by
Andorra's demographic composition: approximately 55\% of residents are
foreign-born labour migrants, and historical data show that net immigration
responds to economic conditions with a lag of one to two years rather than
driving them~\cite{imf2025andorra,worldbank2024}.

\paragraph{Causal chain.}
GDP growth generates labour demand; firms respond by requesting additional
immigration permits; new residents arrive with a lag. The GDP-to-population
elasticity is calibrated at 0.40 at a one-year lag and 0.10 at a two-year lag,
consistent with IMF\,(2025) Andorra Selected Issues estimates and with the
broader labour-migration literature for small European jurisdictions.

\paragraph{Population update rule.}
Let $g^\text{GDPpc}_{t-k}$ denote the growth rate of GDP per capita at lag
$k$, and let $g^\text{nat}=-0.1\%\,\text{yr}^{-1}$ capture the natural
demographic change (deaths exceed births at Andorra's fertility rate of
0.8--1.5). The endogenous population evolves as

\begin{equation}
  P_{t+1} = P_t\!\left(1 + g^\text{nat}
              + \beta_1\, g^\text{GDPpc}_{t}
              + \beta_2\, g^\text{GDPpc}_{t-1}\right)
            - \beta_\text{afford}
              \max\!\left(0,\,\frac{A_t}{100} - 0.30\right)P_t,
  \label{eq:pop_gdp}
\end{equation}
where $\beta_1=0.40$, $\beta_2=0.10$, and $\beta_\text{afford}=0.10$ is the
housing-affordability suppression coefficient described in
Section~\ref{sec:housing_feedback}. The scenario layer specifies a GDP per
capita trajectory (via the TFP growth rate and the tourism driver); population
follows endogenously from Equation~\eqref{eq:pop_gdp} without requiring an
assumed headcount endpoint.

\subsection{Housing-Affordability Migration Feedback}
\label{sec:housing_feedback}

Empirical evidence from small European jurisdictions suggests that housing
affordability materially affects residential location
decisions~\cite{Sa2015,Diamond2016}. Let $A_t$ denote the ratio of median
annual housing cost to median gross household income (expressed as a
percentage), and let $\bar{A}=30\%$ denote the affordability threshold above
which net in-migration is suppressed. The last term of
Equation~\eqref{eq:pop_gdp} implements this feedback: it is zero whenever
housing is affordable and becomes active only once the burden threshold is
breached, ensuring the feedback is one-sided.

Housing prices evolve endogenously as a function of population pressure and
income growth; the resulting affordability ratio $A_t$ feeds back into net
migration as described above. Under all four scenarios the median 2049
affordability burden converges to approximately 30\% of gross household income,
consistent with the threshold at which in-migration ceases to amplify further
housing demand.

\subsection{Tourism Carrying-Capacity Constraint}
\label{sec:tourism_cap}

Following the Tourism Area Life Cycle framework~\cite{Butler1980} and empirical
evidence on small-destination saturation~\cite{Peeters2019}, we impose a
logistic ceiling on annual tourism arrivals. Infrastructure-scaled carrying
capacity is defined as
\begin{equation}
  C_t = C_0 \left(\frac{B_t}{B_0}\right)^{0.4}
              \left(\frac{P_t}{P_0}\right)^{0.2},
  \label{eq:tour_cap}
\end{equation}
where $C_0 = 9{,}646{,}656 / 0.85 \approx 11.35\,\text{M}$ visitor-nights is
the implied capacity at the 2024 base year (with base businesses $B_0=10{,}645$
and base population $P_0=87{,}097$). The exponents reflect the relative roles of
commercial infrastructure (0.4) and resident hosting capacity (0.2) in
determining the absorption limit.

When utilisation $u_t = T_t / C_t$ exceeds the critical threshold
$\bar{u}=0.85$, arrivals are curtailed via an exponential friction term:
\begin{equation}
  T_{t+1}^* =
  \begin{cases}
    T_{t+1} & \text{if } u_t \leq 0.85,\\[4pt]
    C_t \cdot 0.85\;+\;\bigl(T_{t+1} - C_t \cdot 0.85\bigr)
        \exp\!\bigl(-3\,(u_t - 0.85)\bigr) & \text{if } u_t > 0.85.
  \end{cases}
  \label{eq:tour_friction}
\end{equation}
The friction coefficient of 3 implies that each percentage-point of excess
utilisation above the threshold attenuates marginal arrivals by approximately
$3\%$; at $u_t=1.0$ the effective ceiling is approached asymptotically.

\subsection{Scenario Engine Dynamics}
\label{subsec:scenario_engine}

The scenario engine advances the state $S_t \rightarrow S_{t+1}$ using
deterministic equations and bounded constraints. We define the growth operator
\begin{equation}
g(x_t)=\frac{x_t-x_{t-1}}{\max(x_{t-1},\epsilon)},
\end{equation}
and a clamp function $clamp(x,a,b)=\min(\max(x,a),b)$.

\paragraph{Exogenous drivers vs endogenous KPIs.}
In the updated framework, GDP per capita trajectory (set by the TFP growth rate
and the tourism driver) is the primary exogenous scenario lever. Population is
endogenous, following from the GDP-to-labour-immigration chain described in
Section~\ref{sec:gdp_pop_model}. All remaining KPIs evolve endogenously through
shared update rules, which enforces cross-domain coupling (e.g., population
affects housing, housing affects affordability, affordability feeds back into
migration pressure).

\paragraph{Core state-update template.}
For each endogenous KPI $k_t$, the engine applies:
\begin{equation}
k_{t+1} = clamp\Big(k_t + F_k(S_t, D_{t+1};\theta_k),\; \underline{k},\overline{k}\Big),
\end{equation}
where $D_{t+1}$ are scenario drivers at year $t{+}1$ and $\theta_k$ are
calibrated parameters.

\paragraph{Examples of coupled update forms.}

\textit{(1) Stock--flow dynamics for built environment.}
\begin{equation}
Stock_{t+1} = \max(\underline{Stock},\; Stock_t + Inflow_t - \delta\,Stock_t),
\end{equation}
where $Inflow_t$ is a nonnegative function of demand pressure and constraints,
and $\delta$ is a depreciation/demolition rate.

\textit{(2) Price/pressure proxies from scarcity and demand.}
\begin{equation}
Price_{t+1} = H\Big(Income_{t+1},\; \frac{Stock_{t+1}}{Demand_{t+1}},\; Rate_t,\; Driver_{t+1}\Big).
\end{equation}

\textit{(3) Burden-style affordability metrics.}
\begin{equation}
Burden_{t+1}=clamp\left(\frac{Cost_{t+1}/12}{\max(Salary_{t+1},\epsilon)}\cdot 100,\; 0,\; \overline{b}\right).
\end{equation}

\textit{(4) Environmental accounting linked to socio-economic drivers.}
\begin{equation}
Envpc_{t+1}=Envpc_0 \cdot \left(\frac{GDPpc_{t+1}}{GDPpc_0}\right)^{\eta}
\cdot \eta_{dens}(Density_{t+1})\cdot \eta_{driver}(Driver_{t+1}),
\quad
Env_{t+1}=Envpc_{t+1}\cdot Pop_{t+1}.
\end{equation}

\subsection{Derived Infrastructure Calculator}
\label{subsec:infrastructure_calc}

Infrastructure requirements are computed deterministically from the current-year
state, primarily from population and per-capita/service ratios.

\paragraph{Electricity demand and capacity.}
\begin{equation}
E_t = Pop_t \cdot e_{pc}, \qquad C_t = \frac{E_t}{CF \cdot H},
\end{equation}
where $e_{pc}$ is per-capita annual electricity use, $CF$ is a capacity factor,
and $H=8760$ hours/year. A renewable share $r_t$ partitions required capacity:
\begin{equation}
C^{RE}_t = r_t \cdot C_t, \qquad C^{NRE}_t = (1-r_t)\cdot C_t.
\end{equation}

\paragraph{Water demand and security index.}
\begin{equation}
W^{hh}_t = \frac{Pop_t \cdot w_{pc} \cdot 365}{1000}, \qquad
W^{tot}_t = \frac{W^{hh}_t}{\alpha}, \qquad
I^W_t = \frac{Y_{safe}}{W^{tot}_t},
\end{equation}
where $w_{pc}$ is liters/person/day, $\alpha$ is the household share of total
demand, and $Y_{safe}$ is a safe-yield proxy.

\paragraph{Health and education capacity.}
\begin{equation}
Beds^{req}_t = \frac{b}{1000}\,Pop_t, \qquad
Stu_t = \left(\frac{Stu_0}{Pop_0}\right)Pop_t, \qquad
Class_t = \frac{Stu_t}{s_{class}}, \qquad Schools_t = \frac{Stu_t}{s_{school}}.
\end{equation}

\subsection{Reproducibility and Updateability}
\label{subsec:reproducibility}

The methodology is designed to be re-runnable when new statistics arrive. The
KPI dictionary, unit conversions, harmonization rules, scenario update
equations, and spatial allocation scripts are treated as versionable components.
This ensures that scenario outputs remain auditable artifacts of (i) a specific
data release and (ii) an explicit set of scenario drivers and parameters, rather
than an opaque one-off spreadsheet process.

% ============================================================
\section{Experiments}
\label{sec:experiments}
% ============================================================

This section describes the scenario suite implemented in this study and the
experimental protocol used to generate outputs from the Andorra Clone Scenario
Engine over the 2024--2049 horizon. Consistent with the MIT Media Lab City
Science approach, the experiments are framed as comparative decision-support
exercises designed to make trade-offs legible across alternative futures, rather
than as a search for a single optimal forecast~\cite{orii2020methodology}.

\subsection{Experimental Protocol}
\label{subsec:exp_protocol}

All experiments share a common baseline initialization year (2024) and
simulation horizon $T=2049$. Each run proceeds in yearly steps. At each step,
the scenario layer provides (i) the GDP per capita growth target for the year
(the primary exogenous driver) and (ii) scenario-specific constraints or policy
regimes. The engine then updates the KPI state vector $S_t \rightarrow S_{t+1}$
using the deterministic transition rules defined in
Section~\ref{sec:methodology}. Population evolves endogenously from the GDP
trajectory via Equation~\eqref{eq:pop_gdp}.

\paragraph{Controlled components.}
Across scenarios, we hold constant: (i) the KPI dictionary, (ii) the baseline
initialization procedure, (iii) harmonization and clamping rules, and (iv) the
infrastructure calculator ratios. This ensures differences in outputs arise from scenario
assumptions rather than changes in measurement.

\paragraph{Recorded outputs.}
For each run, the system records the full KPI trajectories $\{S_t\}_{t=0}^{T}$ and
derived infrastructure requirements $\{I_t\}_{t=0}^{T}$.

\subsection{Scenario Construction}
\label{subsec:scenario_construction}

Each scenario is defined by two elements:
\begin{enumerate}
    \item \textbf{GDP per capita growth trajectory:} a specified TFP growth rate
          and tourism driver that together determine the GDPpc path, from which
          population follows endogenously.
    \item \textbf{Constraints / policy regime:} hard or soft rules that modify
          how endogenous variables respond (e.g., spatial expansion restrictions
          in Density, or accelerated affordability brake in Continuity).
\end{enumerate}

\subsection{Scenario Suite Implemented}
\label{subsec:scenario_suite}

We implement four scenarios spanning a broad design space of planning stress
tests. The 2049 population endpoints implied by the GDP trajectories are:
Overgrowth~179{,}021; Continuity~117{,}666; Density~104{,}914;
Degrowth~73{,}172.

\paragraph{Continuity (trend-following baseline).}
A continuation pathway in which the TFP growth rate and tourism driver follow
recent historical trends. Functions as the reference case.

\paragraph{Overgrowth (high-growth stress test).}
A high-growth pathway with elevated TFP ($g_\text{TFP}=6.0\%$) and sustained
tourism expansion, implying strong labour demand and large net immigration.
Designed to test upper-bound pressures on housing, infrastructure, and
environmental proxies. Read as a stress scenario, not a normative target.

\paragraph{Degrowth (contraction / low-growth stress test).}
A low-growth pathway testing downward pressures on demand and service
requirements. Tourism contracts and TFP stagnates, reducing labour demand and
triggering net emigration in the low-accessibility periphery.

\paragraph{Density (compact development regime).}
A compact-growth pathway that enforces a no-expansion regime on the spatial
footprint: all new residents are accommodated within existing walkable-access
cores regardless of the GDP-driven immigration flow. Tests intensity-driven
effects under constrained footprint expansion.

\subsection{Comparison Metrics}
\label{subsec:exp_metrics}

Scenario comparison is multi-objective. We evaluate scenarios using:
\begin{itemize}
    \item \textbf{Endpoint comparison:} 2049 KPI values and derived
          infrastructure requirements.
    \item \textbf{Trajectory comparison:} full time-series behavior to identify
          transient effects, nonlinearities, and constraint activation.
    \item \textbf{Normalized composite views:} min--max normalization for
          dashboard comparisons (with inversion for ``lower-is-better'' metrics).
\end{itemize}

% ============================================================
\section{Results}
\label{sec:results}
% ============================================================

\subsection{Baseline (2024) and horizon-year comparison (2049)}
\label{subsec:baseline_and_2049}

All scenario runs start from the same baseline state $S_{2024}$. The 2024
initialization values are reported in the dashboard and serve as the shared
reference for all scenario comparisons.

\begin{table}[htbp]
\centering
\caption{Key 2049 indicators under four scenarios (calibrated model, median
         run). Uncertainty bands (P10--P90) from $N=500$ Monte Carlo draws are
         shown for GDP per capita. Population values are the endogenous model
         outcomes driven by each scenario's GDP trajectory.}
\label{tab:scenarios_2049}
\small
\begin{threeparttable}
\begin{tabularx}{\textwidth}{lXXXX}
\toprule
\textbf{Indicator} & \textbf{Continuity} & \textbf{Overgrowth}
                   & \textbf{Degrowth}   & \textbf{Density} \\
\midrule
Population
  & 117{,}666 & 179{,}021 & 73{,}172 & 104{,}914 \\
Foreign-born share (\%)
  & $\sim$62 & $\sim$77 & $\sim$51 & $\sim$60 \\
Housing afford.\ (\% income)
  & $\approx$30 & $\approx$30 & $<$30 & $\approx$30 \\
Tourism arrivals (M)\tnote{$\dagger$}
  & $\sim$12 & $\sim$17 & $\sim$5.5 & $\sim$12 \\
CO$_2$ per capita (t)
  & moderate & lower & higher & moderate \\
Renewable share
  & $\sim$96\% & $\sim$97\% & $\sim$95\% & $\sim$96\% \\
Water security index
  & declining & low & maintained & declining \\
Hospital beds gap
  & $+57$ & $+230$ & $-93$ & $+208$ \\
\bottomrule
\end{tabularx}
\begin{tablenotes}
  \item[$\dagger$] Tourism ceiling active in Overgrowth and Density;
        unconstrained projections would exceed infrastructure capacity.
\end{tablenotes}
\end{threeparttable}
\end{table}

\begin{table*}[t]
\centering
\small
\caption{Horizon-year (2049) outputs across scenarios. Infrastructure values are
derived from the state vector using fixed service ratios. ``Beds gap'' is
expressed relative to baseline hospital-bed capacity. Population values are
endogenous model outputs driven by the GDP trajectory of each scenario.}
\label{tab:results_2049}
\begin{tabular}{p{2.5cm} p{2.0cm} p{2.0cm} p{2.0cm} p{2.0cm}}
\hline
\textbf{Indicator (2049)} & \textbf{Continuity} & \textbf{Overgrowth} & \textbf{Degrowth} & \textbf{Density} \\
\hline
Population $P$ & 117{,}666 & 179{,}021 & 73{,}172 & 104{,}914 \\
Foreign-born (\%) & 62 & 77 & 51 & 60 \\
Electricity demand (kWh/yr) & $353\times10^6$ & $537\times10^6$ & $219\times10^6$ & $315\times10^6$ \\
Beds gap & $+57$ & $+230$ & $-93$ & $+208$ \\
\hline
Access [0,1] & 0.938 & 0.980 & 0.921 & 0.935 \\
\hline
NatCov & $\sim$0.91 & $\sim$0.86 & $\sim$0.93 & $\sim$0.93 \\
WaterSec & declining & low & maintained & declining \\
AQI & moderate & elevated & low & moderate \\
\hline
\end{tabular}
\end{table*}

\begin{figure}[H]
    \centering
    \includegraphics[width=1\linewidth]{scenario_radar_comparison_2035 (1)}
    \caption{Cross-scenario comparison: radar plot across eight indicators
    (min--max normalized). CO$_2$ and AQI are inverted so higher scores reflect
    better environmental outcomes.}
    \label{fig:scenario_radar_2049}
\end{figure}

\begin{figure}[H]
    \centering
\begin{subfigure}[t]{0.5\linewidth}
    \centering
    \includegraphics[width=\linewidth]{population_scenarios_2025_2035 (1)}
    \caption{Population $P(t)$ --- endogenous, GDP-driven}
    \label{fig:population}
\end{subfigure}\hfill
\begin{subfigure}[t]{0.5\linewidth}
    \centering
    \includegraphics[width=\linewidth]{gdp_per_capita_scenarios_2025_2035 (1)}
    \caption{GDP per capita $GDPpc(t)$ --- primary exogenous lever}
    \label{fig:gdp-per-capita}
\end{subfigure}

\vspace{0.8em}

\begin{subfigure}[t]{0.5\linewidth}
    \centering
    \includegraphics[width=\linewidth]{co2_total_scenarios_2025_2035 (1)}
    \caption{Total CO$_2$ Emissions}
    \label{fig:co2}
\end{subfigure}\hfill
\begin{subfigure}[t]{0.5\linewidth}
    \centering
    \includegraphics[width=\linewidth]{forest_coverage_scenarios_2025_2035 (1)}
    \caption{Natural Coverage}
    \label{fig:natcov}
\end{subfigure}

\caption{Time-series trajectories (2014--2049). Historical values (2014--2024)
provide calibration context; projections (2025--2049) are scenario outputs.
Population trajectories are endogenous outputs of the GDP-driven immigration
model rather than imposed inputs.}
\label{fig:timeseries-2x2}
\end{figure}

\subsection{Trade-offs and output consistency}
\label{subsec:tradeoffs_and_consistency}

Because the system couples socio-economic, environmental, and infrastructure
indicators, there is no universal ``best'' scenario; desirability depends on
planning priorities. The results reveal four dominant trade-offs:

\begin{itemize}
    \item \textbf{Economic scale vs.\ capacity expansion:} Higher GDP growth
          drives larger immigration flows and higher population, which in turn
          raises electricity demand and service requirements (beds, schools).
          In the highest-growth pathway, environmental proxies also deteriorate
          (higher $AQI$ and lower $NatCov$), indicating coupled
          infrastructure--environment stress.
    \item \textbf{Business-as-usual does not remove resource stress:} The
          trend-following pathway produces moderate growth, yet water security
          declines relative to baseline and total emissions rise, implying
          continuity does not eliminate key constraints.
    \item \textbf{Environmental relief via contraction:} Degrowth improves
          $NatCov$, $AQI$, and water security while increasing $WLB$, but
          reduces GDP-driven labour demand and implies downsizing of service
          capacity (negative bed gap and fewer schools).
    \item \textbf{Compact growth shifts pressure into intensity:} Under a
          no-expansion spatial regime, natural cover can be preserved while
          GDP-driven immigration continues, but pressures reappear as
          intensity-driven outcomes (lower water security and higher required
          services per fixed networks), highlighting the tension between land
          preservation and per-capita provision.
\end{itemize}

% ============================================================
\section{Discussion}
\label{sec:discussion}
% ============================================================

\subsection{Are the results ``good'' or ``bad''? (fitness-for-purpose)}

Results should be judged against the system's goal: producing decision-relevant,
cross-domain-consistent scenario trajectories for a small-country planning
context, not a single ``best forecast.'' By this criterion, results are
\textbf{good} in three ways: (i) outputs are \textbf{coupled}, so changes in
the GDP driver propagate coherently into population, infrastructure requirements,
and environmental pressures via the labour-immigration chain; (ii) scenarios
\textbf{separate planning regimes}, surfacing qualitatively different stress
patterns; and (iii) trade-offs remain \textbf{interpretable}, enabling
comparative stress testing.

However, results are \textbf{not yet ``good'' as predictive truth} for all
indicators: several outcomes are proxy-based and require further calibration and
validation before certain quantities can be treated as absolute rather than
comparative. The extended 2049 horizon also amplifies compounding uncertainties,
particularly in the GDP-to-population elasticity and the TFP growth rate.

\subsection{Connecting findings to the research question}

The research question asks whether a transparent, reproducible scenario engine
can generate decision-relevant futures for a small country under uncertainty
while maintaining cross-domain consistency. The results support this: outputs
appear as \emph{bundled futures} in which socio-economic indicators,
environmental signals, and infrastructure requirements co-evolve from a shared
state update, making trade-offs explicit. The inverted causal chain---GDP drives
population rather than vice versa---removes the need to impose arbitrary
population headcount assumptions and grounds the demographic projections in
observed labour-market dynamics.

\subsection{Credibility and cautions}

\paragraph{More credible (structural outputs).}
The most defensible results follow directly from accounting and deterministic
scaling: infrastructure requirements from population and service ratios,
land-pressure measures tied to buildable constraints, and directionally
consistent coupled dynamics. The GDP-to-population elasticity is calibrated
against observed IMF and World Bank data.

\paragraph{Use with caution (proxy outcomes).}
Wellbeing indices and selected environmental proxies should be treated as
\emph{comparative signals} rather than point predictions. The GDP-to-population
elasticity, while consistent with published estimates, is estimated from a short
and partially coincident time series; structural breaks (e.g., post-COVID
labour-market reorganization) could shift it materially. The 2049 horizon
amplifies all compounding errors relative to the original 2035 projection range.

\subsection{Comparison with alternative methodologies}

Time-series forecasting can be statistically grounded for individual variables
but often lacks cross-domain coupling. System dynamics captures feedbacks but
can become complex and calibration-heavy. ML can perform well with abundant
stable data, but small-country multi-domain settings are typically data-limited
and require explicit feasibility constraints. The scenario engine trades
predictive optimality for transparency, updateability, and cross-domain
consistency. The inversion of the population--GDP causal direction is a
structural change that aligns the model more closely with Andorra's actual
demographic drivers and reduces the arbitrariness of scenario headcount
assumptions.

\subsection{Planning implications}

The results suggest that planning choices should be evaluated as \emph{bundles}
rather than single KPI targets: higher economic growth drives larger immigration
flows and higher service requirements; contraction relieves some pressures but
implies difficult service right-sizing; compact growth can preserve land take
but concentrates pressure into intensity-driven constraints (affordability and
per-capita capacity). The housing-affordability feedback mechanism implies that
the rent-to-income ratio is projected to approach 30\% across most pathways,
activating the in-migration brake. This threshold should be treated as an
early-warning indicator for housing policy. On tourism, the capacity model
implies that sustained growth above 11--12 million visitor-nights per year will
require commensurate investment in commercial infrastructure or active demand
management.

\subsection{Limitations}

\begin{enumerate}
    \item \textbf{Scenario forcing vs.\ causality:} GDP trajectories are steered
          toward scenario-defined growth rates; outputs are internally consistent
          scenarios, not causal forecasts.
    \item \textbf{GDP-to-population elasticity uncertainty:} the elasticity
          coefficients ($\beta_1=0.40$, $\beta_2=0.10$) are calibrated against
          a relatively short historical series. Post-COVID structural changes in
          Andorra's labour market could shift these values. A $\pm50\%$
          sensitivity test on $\beta_1$ shifts the 2049 Continuity population
          by approximately $\pm18{,}000$ persons.
    \item \textbf{Housing market dynamics:} the model does not distinguish
          owner-occupied from rental tenures, nor does it capture speculative
          dynamics or the role of non-resident investors. Supply-side constraints
          are held constant.
    \item \textbf{Proxy indices:} wellbeing and resource-security indices remain
          simplified and should be treated as comparative signals until
          parameterized with domain-specific data.
    \item \textbf{Tourism capacity parameterization:} the capacity function is
          calibrated on a single cross-sectional observation (2024) and the
          infrastructure elasticities are assumed rather than empirically
          estimated from panel data. Seasonal distribution of arrivals, visitor
          heterogeneity, and the quality dimension of tourism are not represented.
    \item \textbf{Scope of uncertainty quantification:} Monte Carlo propagation
          covers parametric uncertainty in $g_\text{TFP}$, $\gamma_2$, and
          $a_1$ only; structural uncertainty and deep uncertainty over exogenous
          shocks lie outside the present scope. The extended 2049 horizon
          amplifies these gaps relative to the original 2035 projections.
\end{enumerate}

\subsection{Next steps (toward V2 validation)}

Beyond improving calibration of the KPI engine, the main V1$\rightarrow$V2
upgrade is a \textbf{humanized AI agent layer} inside Andorra Clone to simulate
social response under alternative futures. The goal is a population of agents
representing people and groups in Andorra who perceive, react, and decide in
locally coherent ways, adding evaluation signals beyond macro KPIs such as
perceived wellbeing and social
acceptance~\cite{adornetto2025generative,garcia2025simulating,carrero2024integrating,li2025hugagent}.

\paragraph{Additional V2 upgrades.}
\begin{itemize}
    \item Empirical re-estimation of the GDP-to-population elasticity using
          panel data from comparable Alpine micro-states (Liechtenstein, San
          Marino, Monaco) to reduce the reliance on a single country time series.
    \item Formal KPI metadata schema and automated tests (definitions, units,
          temporal resolution, provenance, bounds, transformations) to prevent
          scale drift across releases.
    \item Empirical estimation of tourism infrastructure elasticities from panel
          data across comparable Alpine microstates.
    \item Coupling the environmental module to a regional climate downscaling
          product (e.g., Pyrenees-CORDEX) to endogenise temperature forcing.
\end{itemize}

% ============================================================
\section{Conclusion}
% ============================================================

This paper presented the Andorra Clone Scenario Engine, a transparent,
data-driven framework that transforms heterogeneous public and governmental
indicators into comparable long-horizon trajectories to 2049. Three
methodological advances distinguish this version. First, a
\textbf{causal inversion of the population--economy relationship}: rather than
treating population as an exogenous driver, the engine derives population
endogenously from GDP growth via a labour-immigration elasticity
($\beta_1=0.40$ at one-year lag, $\beta_2=0.10$ at two-year lag), grounded in
IMF\,(2025) Andorra Selected Issues and consistent with labour-market dynamics
in small European jurisdictions. This removes the need for arbitrary headcount
assumptions and produces demographically coherent scenario trajectories.
Second, key macroeconomic parameters are calibrated against the 2010--2024
historical record via Nelder--Mead optimisation (RMSE\,=\,1.77\%), anchoring
all projections in observed dynamics.

Across scenarios, the results show clear system-level trade-offs. Under the
GDP-driven demographic model, the 2049 population ranges from 73{,}172
(Degrowth) to 179{,}021 (Overgrowth), with Continuity at 117{,}666 and Density
at 104{,}914. Expansion-driven growth tends to increase economic throughput
while amplifying housing pressure, infrastructure load, and environmental stress.
Degrowth reduces demand-side pressures and improves several environmental
proxies, but constrains economic output and implies right-sizing of service
infrastructure. The Density scenario separates \emph{more people} from
\emph{more land use}: maintaining a fixed built footprint can preserve nature
coverage, but concentrates stress in affordability and service provision within
existing walkable cores.

The housing-affordability feedback mechanism indicates that the rent-to-income
ratio approaches 30\% across most pathways, activating the in-migration brake
and suggesting that housing affordability is a binding planning constraint under
all scenarios. Tourism arrivals are capped by the logistic carrying-capacity
model near 17 million visitor-nights in Overgrowth, roughly 30\% below
unconstrained extrapolations, highlighting the importance of active demand
management at high-growth levels.

The next major extension---a humanized AI agent layer---will allow evaluation
not only of macro feasibility but of social acceptance and lived experience
across Andorra's diverse resident communities. In line with the MIT Media Lab
City Science vision, the framework is designed throughout as a transparent,
human-in-the-loop decision-support system for exploring trade-offs across
plausible futures rather than optimizing a single deterministic outcome.

% bibliography at end
\bibliographystyle{plain}
\bibliography{references}

\end{document}
